\documentclass[letterpaper,journal]{IEEEtran}

\usepackage{amsmath,amsfonts,amssymb,bm}
\usepackage{algorithmic}
\usepackage{algorithm}
\usepackage{array}
\usepackage{textcomp}
\usepackage{stfloats}
\usepackage{url}
\usepackage{verbatim}
\usepackage{graphicx}
\usepackage{cite}
\usepackage{booktabs}
\usepackage{multirow}
\usepackage{xcolor}

\hyphenation{op-tical net-works semi-conduc-tor IEEE-Xplore}

% -----------------------------------------------------------------------------
% Draft controls. Set \draftnotesfalse before submission and resolve every note.
% -----------------------------------------------------------------------------
\newif\ifdraftnotes
\draftnotestrue
\newcommand{\draftnote}[1]{\ifdraftnotes\textcolor{red}{\textbf{[TODO: #1]}}\fi}
\newcommand{\Bern}{\operatorname{Bernoulli}}
\newcommand{\Normal}{\mathcal{N}}
\newcommand{\PG}{\operatorname{PG}}
\newcommand{\KL}{D_{\mathrm{KL}}}
\newcommand{\sigmoid}{\operatorname{sigmoid}}
\newcommand{\E}{\mathbb{E}}
\newcommand{\Var}{\operatorname{Var}}
\newcommand{\given}{\,\middle|\,}

% The paths below point to the executable snapshots used for the preliminary
% tables in this draft. Replace them with frozen, confirmatory artifacts.
\graphicspath{{../../../outputs/fully_bayesian/20260729_191716/}}
\newcommand{\resultgraphic}[2]{\includegraphics[width=#2]{#1}}

\begin{document}

\title{Few-Feedback Personalization of Binary Ride Preference over Road Bumps Using Hierarchical Bayesian Learning}

% T-IV uses single-blind review. Replace this placeholder with the real author
% list and add affiliations/biographies in the journal-prescribed locations.
\author{Anonymous Authors}

\maketitle

\begin{abstract}
Ride preference varies across occupants, whereas most data-driven ride evaluators produce a population-level score. We present an uncertainty-aware hierarchical Bayesian framework that predicts whether an evaluator judges a road-bump response as good or bad and adapts to a previously unseen evaluator from limited feedback. Interpretable transient and frequency-weighted features are extracted from inertial and body-motion signals. Evaluator-specific logistic coefficients share a population distribution, and group spike-and-slab indicators select useful motion features. P\'{o}lya--Gamma augmentation supports posterior sampling for population learning and personalization. Projection-predictive and cross-fitted conditional-information analyses check feature sufficiency. Evaluation includes a high-feedback protocol and a sparse-feedback protocol with 10--35 labels per unseen evaluator, using discrimination, probability accuracy, log predictive density, and posterior uncertainty. In the current executable snapshot, personalization changes high-feedback-target AUROC from 0.843 to 0.846 and Brier score from 0.211 to 0.190. Across five sparse-feedback targets, macro-average Brier score improves from 0.244 to 0.205, although individual AUROC estimates remain highly uncertain. Replacing one body-motion channel with a driver-held-out recurrent reconstruction is evaluated only as a robustness study. These preliminary results motivate a confirmatory evaluator-held-out study with repeated sampling and stronger baselines.
\end{abstract}

\begin{IEEEkeywords}
Bayesian learning, human factors, intelligent vehicles, ride preference.
\end{IEEEkeywords}

\section{Introduction}
\label{sec:introduction}

Ride quality is simultaneously a vehicle-dynamics property and a human judgment.
Objective quantities such as acceleration peaks, vibration dose value, and frequency-weighted
root-mean-square acceleration are useful engineering descriptors, but they do not uniquely
determine whether a particular occupant likes a particular transient response. The recent review
by Guastadisegni \emph{et al.} emphasizes both the value of correlating objective measurements
with subjective attributes and the need for nonlinear, multivariable models of human ride
sensitivity~\cite{guastadisegni2024ride}. This need becomes more consequential in intelligent
vehicles, where control policies and software-defined chassis settings can in principle be adapted
to the current occupant rather than tuned only for an average customer.

Three statistical difficulties arise in this setting. First, preferences differ across people, so
pooling all labels into one classifier can erase meaningful heterogeneity. Second, a new occupant
usually supplies only a few labels; fitting an independent model is then unstable. Third, an
intelligent vehicle should distinguish uncertainty about the occupant from irreducible ambiguity in
the feedback. A point estimate alone cannot indicate whether the system has observed enough
feedback to personalize safely.

We address these difficulties with a hierarchical Bayesian binary-response model. Each evaluator
has an individual coefficient vector, while the population mean and covariance allow partial
pooling across evaluators. A previously unseen evaluator begins from the population posterior and
is updated sequentially using their own good/bad feedback. Exact-zero group spike-and-slab
variables, projection-predictive search, and conditional-information analysis are used to study
which physically interpretable motion features support prediction. The framework targets a
human-factors capability within the scope of intelligent vehicles: learning an occupant-specific
ride evaluation from in-vehicle sensing and sparse interaction.

The contributions of this work are as follows.

\begin{itemize}
  \item We formulate road-bump good/bad preference as a hierarchical Bayesian logistic model that
  transfers population knowledge to a new evaluator and exposes the full posterior predictive
  distribution during sequential personalization.

  \item We combine exact-zero feature inclusion with two predictive diagnostics---projection to
  sparse submodels and cross-fitted conditional-information gain---to separate population-level
  relevance, predictive redundancy, and evaluator heterogeneity among ride-response features.

  \item We define two evaluator-held-out protocols: a high-feedback target experiment that traces
  adaptation over many labels, and a sparse-feedback experiment that stresses 10--35-label target
  histories. Probability scores and uncertainty intervals are reported alongside ranking metrics.

  \item As a secondary robustness study, we replace the body bounce-rate channel with a causal
  recurrent reconstruction learned without the target evaluator and measure the resulting change
  in preference prediction.
\end{itemize}

The numerical results in this manuscript are an evidence-bound preliminary snapshot of the current
implementation. They are included to make the draft executable and to expose limitations; the
confirmatory study required for submission is specified in Section~\ref{sec:confirmatory}.

\section{Related Work}
\label{sec:related}

\subsection{Objective and Subjective Ride Evaluation}

Whole-body vibration standards define frequency-weighted quantities for assessing human exposure,
and vehicle studies commonly use acceleration, jerk, peak, and vibration-dose descriptors
~\cite{iso2631,elbanhawi2015passenger}. Such quantities are valuable priors for representation
design, but an inertial measurement unit mounted near the vehicle center is not a measurement at a
seat, backrest, or foot contact point. Consequently, the weighted quantities in this work are
\emph{ISO-inspired vehicle-response features}; they are not presented as standards-compliant
occupant exposure values.

Machine-learning studies have related acceleration and participant characteristics to subjective
ride comfort~\cite{cieslak2020ride}, and a neural evaluator has been developed specifically for
speed-bump measurements~\cite{kim2019evaluator}. These studies demonstrate that multichannel
signals can predict subjective judgments, but a single pooled evaluator does not directly solve
online cold-start personalization. Recent sensory-driven work explicitly accounts for differences
in vibration perception through perception groups and bootstrap uncertainty
~\cite{kikuchi2025personalizing}. Our model instead maintains a posterior coefficient distribution
for each individual and adapts it from a shared continuous population distribution.

\subsection{Preference-Aware Intelligent Vehicles}

Preference-aware planning and control have been proposed to make automated-vehicle motion more
consistent with occupant-specific acceleration and jerk limits~\cite{bae2020robocar}. Ran
\emph{et al.} use Bayesian pairwise queries to learn preferred lane-centering trajectories online
~\cite{ran2023online}. These approaches establish the value of explicit preference elicitation.
The present problem differs in both feedback and output: feedback is a pointwise good/bad judgment
after a detected bump, and the learned quantity is a posterior evaluator model that can later serve
as an input to vehicle tuning or control. Closing that control loop is outside the scope of this
paper.

\subsection{Bayesian Computation and Feature Selection}

P\'{o}lya--Gamma augmentation turns the logistic likelihood into a conditionally Gaussian form,
enabling efficient Gibbs updates~\cite{polson2013pg}. Spike-and-slab modeling attaches posterior
inclusion probabilities to candidate variables~\cite{george1993spike}. Because inclusion alone can
be unstable when ride features are correlated, we also use projection-predictive selection, which
seeks a small submodel that preserves a reference model's predictions
~\cite{piironen2020projective}. Finally, conditional-information criteria help distinguish feature
relevance from redundancy~\cite{brown2012conditional}. These tools answer different questions and
are therefore reported together rather than collapsed into one importance ranking.

\section{Problem Formulation}
\label{sec:problem}

Let evaluator $u\in\{1,\ldots,U\}$ provide binary feedback for road-bump episode $i$.
The outcome
\begin{equation}
 y_{ui}=\begin{cases}
 1,&\text{the evaluator reports a good ride response},\\
 0,&\text{the evaluator reports a bad ride response}
 \end{cases}
\end{equation}
is an operational preference label, not a standardized comfort rating. A synchronized sensor window
$\bm{x}_{ui}\in\mathbb{R}^{T\times C}$ is mapped to a standardized feature vector
$\bm{z}_{ui}=\phi(\bm{x}_{ui})\in\mathbb{R}^{d}$, including an intercept.

The training set contains multiple evaluators with both label classes. At deployment, an unseen
evaluator $\star$ initially has no personal labels. After receiving a chronological context
$\mathcal{D}_{\star,t}=\{(\bm{z}_{\star i},y_{\star i})\}_{i=1}^{t}$, the model must estimate
\begin{equation}
 p_{\star j,t}=p(y_{\star j}=1\mid \bm{z}_{\star j},
          \mathcal{D}_{\star,t},\mathcal{D}_{\mathrm{pop}})
\end{equation}
for held-out episodes $j$. The main questions are (i) how much a population prior helps at $t=0$,
(ii) how quickly personal feedback improves probability estimates, and (iii) how uncertain those
estimates remain when $t$ is small.

\section{Hierarchical Bayesian Preference Model}
\label{sec:model}

\subsection{Population and Evaluator-Specific Coefficients}

For evaluator $u$, let $\widetilde{\bm\theta}_u$ be a dense latent coefficient vector and let
$\bm\gamma$ be a shared binary inclusion vector. The effective coefficients are
$\bm\beta_u=\bm\gamma\odot\widetilde{\bm\theta}_u$. The observation model is
\begin{align}
 y_{ui}\mid\bm\beta_u &\sim \Bern\!\left(\sigmoid
             (\bm z_{ui}^{\mathsf T}\bm\beta_u)\right),
 \label{eq:likelihood}\\
 \widetilde{\bm\theta}_u\mid\bm\mu,\bm\Sigma
       &\sim \Normal(\bm\mu,\bm\Sigma).                 \label{eq:userprior}
\end{align}
Thus, $\bm\mu$ captures a population tendency while $\bm\Sigma$ captures between-evaluator
variation and correlation among preferences. The conjugate population prior is
\begin{align}
 \bm\Sigma &\sim \operatorname{InvWishart}(\nu_0,\bm\Lambda_0),\\
 \bm\mu\mid\bm\Sigma &\sim
       \Normal(\bm m_0,\bm\Sigma/\kappa_0).             \label{eq:niw}
\end{align}
The executable configuration sets $\bm m_0=\bm 0$, $\kappa_0=1$,
$\nu_0=d+2$, and $\bm\Lambda_0=\bm I_d$. These weak scale choices are applied after feature
standardization and must be assessed by prior-predictive sensitivity checks in the confirmatory
study.

\subsection{Exact-Zero Group Spike-and-Slab Prior}

Let $g=1,\ldots,G$ index either a single feature or all features derived from one sensor. The
inclusion prior is
\begin{align}
 \gamma_g\mid\pi &\sim \Bern(\pi), &
 \pi &\sim \operatorname{Beta}(a,b),                  \label{eq:spikeslab}
\end{align}
with $a=b=1$. A unit with $\gamma_g=0$ has an exactly zero effective coefficient for every
evaluator; the intercept is always included. Posterior inclusion probability
$\Pr(\gamma_g=1\mid\mathcal D)$ summarizes whether a feature is supported across population
draws. Because correlated features can substitute for one another, this probability is interpreted
together with predictive projection and conditional-information results.

\subsection{P\'{o}lya--Gamma Gibbs Sampling}

For compactness, define $\kappa_{ui}=y_{ui}-1/2$ and
$\bm X_u=\bm Z_u\operatorname{diag}(\bm\gamma)$. Introducing
\begin{equation}
 \omega_{ui}\mid\widetilde{\bm\theta}_u
 \sim \PG\!\left(1,(\bm\gamma\odot\bm z_{ui})^{\mathsf T}
                       \widetilde{\bm\theta}_u\right)
\end{equation}
makes the coefficient conditional Gaussian. With
$\bm\Omega_u=\operatorname{diag}(\bm\omega_u)$, its precision and mean are
\begin{align}
 \bm A_u &= \bm X_u^{\mathsf T}\bm\Omega_u\bm X_u
              +\bm\Sigma^{-1},\\
 \bm m_u &= \bm A_u^{-1}\left(
              \bm X_u^{\mathsf T}\bm\kappa_u
              +\bm\Sigma^{-1}\bm\mu\right),\\
 \widetilde{\bm\theta}_u\mid - &\sim
              \Normal(\bm m_u,\bm A_u^{-1}).           \label{eq:pgconditional}
\end{align}
The sampler alternates the P\'{o}lya--Gamma variables, group indicators, inclusion rate,
evaluator coefficients, and normal--inverse-Wishart population parameters. The current snapshot
uses 500 burn-in iterations and retains 1500 draws without thinning.

\subsection{Cold Start and Sequential Personalization}
\label{sec:personalization}

For posterior draw $m$, a new evaluator is initialized by
\begin{align}
 \widetilde{\bm\theta}_{\star}^{(m)} &\sim
  \Normal(\bm\mu^{(m)},\bm\Sigma^{(m)}),\\
 \bm\beta_{\star}^{(m)} &=
  \bm\gamma^{(m)}\odot\widetilde{\bm\theta}_{\star}^{(m)}.
\end{align}
This is the cold-start population prior. After context labels arrive, eight conditional
P\'{o}lya--Gamma sweeps update $\widetilde{\bm\theta}_{\star}^{(m)}$ for every population draw
while retaining the corresponding $\bm\gamma^{(m)}$. Predictions integrate over population and
personal uncertainty:
\begin{equation}
 \widehat p_{\star j,t}=\frac{1}{M}\sum_{m=1}^{M}
 \sigmoid\!\left(\bm z_{\star j}^{\mathsf T}
                  \bm\beta_{\star,t}^{(m)}\right).      \label{eq:prediction}
\end{equation}

For an episode $j$, we summarize epistemic and Bernoulli uncertainty as
\begin{align}
 U_{\mathrm{epi},j} &= \Var_m(p_j^{(m)}),\\
 U_{\mathrm{ale},j} &= \E_m[p_j^{(m)}(1-p_j^{(m)})].   \label{eq:uncertainty}
\end{align}
The first quantity can contract as personal evidence accumulates; the second remains high for
intrinsically ambiguous posterior probabilities.

\begin{algorithm}[t]
\caption{Sequential personalization of an unseen evaluator}
\label{alg:personalization}
\begin{algorithmic}[1]
\STATE Draw $\widetilde{\bm\theta}_{\star}^{(m)}$ from each population draw
      $(\bm\mu^{(m)},\bm\Sigma^{(m)},\bm\gamma^{(m)})$
\STATE Predict held-out episodes using the cold-start posterior
\FOR{$t\in\{10,20,\ldots,n_{\mathrm{ctx}}\}$, including the final size}
  \STATE Update all $\widetilde{\bm\theta}_{\star}^{(m)}$ using context $1{:}t$
  \STATE Compute posterior predictive probabilities and uncertainty
\ENDFOR
\STATE Report the final context result; retain the best intermediate result only as an oracle diagnostic
\end{algorithmic}
\end{algorithm}

\section{Ride-Response Representation}
\label{sec:features}

\subsection{Windowing and Preprocessing}

The logger records a nominal 10-s episode at 100~Hz, spanning approximately 5~s before and 5~s
after the detected bump event. The classifier uses only the interval from 2~s before to 2~s after
that event. A second-order
10-Hz low-pass filter is applied before downsampling by two, yielding a 50-Hz analysis signal.
Five channels are loaded: 6-D sensor pitch rate and bounce rate, and inertial vertical,
longitudinal, and lateral acceleration. The current manually selected representation uses 12
statistics from bounce rate, vertical acceleration, and longitudinal acceleration; the other loaded
channels remain available for ablations.

All non-intercept features are standardized with means and standard deviations computed from
population-training evaluators only. For a signal $x_1,\ldots,x_T$ at sampling rate $f_s$, the
implemented transient descriptors include
\begin{align}
 \operatorname{P2P}(x) &= \max_t x_t-\min_t x_t,\\
 I_{\mathrm{abs}}(x) &= f_s^{-1}\sum_t |x_t|,\\
 \operatorname{Crest}(x) &=
      \frac{\max_t|x_t|}{\sqrt{T^{-1}\sum_t x_t^2}+10^{-12}},\\
 \operatorname{VDV}(x) &=
      \left(f_s^{-1}\sum_t x_t^4\right)^{1/4}.          \label{eq:features}
\end{align}
Derivatives use a sampling-interval-aware numerical gradient. Frequency-weighted root-mean-square
features multiply the discrete Fourier transform by the magnitude shape of the ISO $W_k$ vertical
or $W_d$ horizontal weighting before inverse transformation and root-mean-square aggregation.
Because the implementation does not reproduce a calibrated seat-location measurement chain, we do
not interpret the resulting value as an ISO exposure measurement.

\begin{table}[t]
\caption{Current Manual Feature Set}
\label{tab:features}
\centering
\setlength{\tabcolsep}{3.5pt}
\begin{tabular}{p{0.30\columnwidth}p{0.62\columnwidth}}
\toprule
Signal & Statistics \\
\midrule
6-D bounce rate & Absolute impulse; derivative absolute peak; derivative peak-to-peak; derivative $W_k$-weighted RMS \\
Vertical acceleration & Peak-to-peak; $W_k$-weighted RMS; crest factor; vibration dose value \\
Longitudinal acceleration & Standard deviation; $W_d$-weighted RMS; absolute impulse; crest factor \\
\bottomrule
\end{tabular}
\end{table}

\subsection{Complementary Feature Analyses}

Projection-predictive selection treats the fitted hierarchical model as a reference. For candidate
subset $S$ and reference probability $p_{ui}^{(m)}$, it solves
\begin{equation}
\begin{aligned}
 \widehat{\bm b}_{S}^{(m,u)}
 &=\underset{\bm b}{\operatorname{arg\,min}}\;\frac{1}{n_u}\sum_i
 \KL\!\left[\Bern(p_{ui}^{(m)})\,\middle\|\right.\\
 &\hspace{25mm}\left.\Bern\{\sigmoid(\bm z_{ui,S}^{\mathsf T}\bm b)\}\right].
\end{aligned}
 \label{eq:projection}
\end{equation}
Forward search greedily adds a sensor or feature. The current stopping rule selects the first
subset recovering 95\% of the reduction from the intercept-only KL divergence to the full model,
using 12 evenly spaced posterior draws.

The conditional-information diagnostic uses five within-evaluator stratified folds. At a forward
step with selected set $S$, candidate $A$ is scored by the cross-fitted log-score gain
\begin{equation}
 \widehat I(Y;Z_A\mid Z_S,U)=\frac{1}{U}\sum_u\frac{1}{n_u}
 \sum_i\left[\ell_{ui}^{S\cup A,-k(i)}-\ell_{ui}^{S,-k(i)}\right].
 \label{eq:cmi}
\end{equation}
The auxiliary logistic models contain evaluator-specific intercepts and slopes. Equation
~\eqref{eq:cmi} is a model-based estimate of conditional information, not a nonparametric exact
mutual-information estimator.

\section{Secondary Missing-Signal Robustness Study}
\label{sec:reconstruction}

The main preference model assumes that 6-D body bounce rate is available. To test sensitivity to
that assumption, a causal sequence-to-sequence long short-term memory network
~\cite{hochreiter1997lstm} reconstructs bounce, roll, and pitch rates from 13 deployable channels:
five inertial signals, four wheel speeds, and four motor-torque quantities. The network has one
96-unit recurrent layer and a linear three-channel output (42,915 trainable parameters). It is
trained for 45 epochs using mean squared error, batch size 32, and learning rate 0.005. All
normalization statistics and network weights exclude every target evaluator.

The preference pipeline is then refit after replacing measured 6-D channels in memory with the
reconstruction. Although three channels are reconstructed, only reconstructed bounce rate affects
the current 12-feature classifier in Table~\ref{tab:features}. This experiment therefore measures
robustness to a virtual bounce-rate sensor; it is not presented as a new reconstruction method or
as uncertainty-aware measurement-error inference.

\section{Experimental Protocol}
\label{sec:experiment}

\subsection{Data and Evaluator Splits}

The repository catalog contains 2891 real-vehicle episodes from 18 evaluators, of which 2682 have
binary labels (1450 good and 1232 bad). The current executable experiments use 10 evaluators. Names
are replaced by protocol-local identifiers, and no name is used as a model input. Table
~\ref{tab:protocols} separates the two intended use cases.

\begin{table*}[t]
\caption{Evaluator-Held-Out Protocols in the Current Executable Snapshot}
\label{tab:protocols}
\centering
\setlength{\tabcolsep}{5pt}
\begin{tabular}{p{0.17\textwidth}cccp{0.14\textwidth}p{0.24\textwidth}}
\toprule
Protocol & \shortstack{Population\\evaluators} & \shortstack{Population\\labels} & \shortstack{Target\\evaluators} & Target labels & Context/holdout construction \\
\midrule
High-feedback target & 5 & 2325 & 1 & 205 & First 102 target labels form the growing context; the remaining 103 form a fixed holdout. \\
Sparse-feedback targets & 5 & 2500 & 5 & 10--35 each (127 total) & First half of each target history is context (5--17); second half is holdout (5--18, 64 total). \\
\bottomrule
\end{tabular}
\end{table*}

The population model never uses labels from a target evaluator. Feature scaling and feature
selection are also fit only on population evaluators. Context and holdout are split in recorded
order rather than randomly, making the test closer to online adaptation but leaving possible
session, vehicle, location, and repeated-bump dependence to be addressed by grouped sensitivity
analyses.

\draftnote{Before submission, state the exact participant recruitment criteria, demographics,
vehicle models/settings, road/bump protocol, feedback interface and timing, compensation,
informed-consent procedure, and the formal institutional ethics determination with approval or
exemption identifier. Do not claim approval until it has been issued. Finalize and justify the
eligibility rule for the eight cataloged evaluators not used in these protocols.}

\subsection{Compared Prediction Stages}

For each target evaluator we report the following stages.

\begin{itemize}
  \item \emph{Prior}: cold-start prediction from the population posterior with zero target labels.
  \item \emph{Final}: prediction after all labels in the allowed context have been assimilated.
  \item \emph{Peak}: the context size with the highest holdout area under the receiver-operating-
  characteristic curve. This uses holdout labels for selection and is therefore an oracle
  diagnostic, never a deployable result.
\end{itemize}

The raw-channel and reconstructed-channel variants are trained from the same evaluator split. A
complete paper must additionally compare pooled logistic regression, independent per-evaluator
models where feasible, regularized partial-pooling baselines, and competitive nonlinear models
under identical evaluator-held-out splits.

\subsection{Metrics and Uncertainty}

We report the area under the receiver-operating-characteristic curve (AUROC), area under the
precision--recall curve (AUPRC), Brier score~\cite{brier1950verification}, and mean log posterior
predictive density
\begin{equation}
 \operatorname{MLPD}=\frac{1}{n}\sum_i
 \log\left[y_i\widehat p_i+(1-y_i)(1-\widehat p_i)\right].
\end{equation}
Ranking metrics assess discrimination, whereas Brier score and log predictive density penalize
poor probabilities. For AUROC uncertainty, each of 600 replicates samples one posterior draw and
bootstraps holdout episodes, and the 2.5th and 97.5th percentiles form an interval. This combines
posterior and finite-holdout variation but does not correct dependence among episodes. Exact or
stratified refitting can additionally provide pointwise expected log predictive density following
Bayesian predictive-evaluation practice~\cite{vehtari2017loo}.

\draftnote{Seed the new-evaluator random-number generator explicitly, run multiple independent
chains and seeds, report effective sample size and rank-normalized convergence diagnostics, and
replace episode bootstrap intervals with evaluator/session/bump-group-aware intervals where the
data hierarchy permits.}

\section{Preliminary Results}
\label{sec:results}

\subsection{High-Feedback Target}

Table~\ref{tab:highfeedback} reports the saved run in which the unseen target provides 102 context
labels and 103 episodes remain fixed for evaluation. On measured channels, personalization changes
AUROC only from 0.8429 to 0.8460, but improves Brier score from 0.2109 to 0.1901 and mean log
predictive density from $-0.6104$ to $-0.5606$. Thus, in this split, the clearest benefit is better
probability quality rather than a large change in ranking. The joint posterior--bootstrap AUROC
interval is $[0.747,0.917]$ for the final measured-channel model; the cold-start interval is much
wider, $[0.207,0.889]$, because population-posterior draws disagree about the target ranking.

\begin{table}[t]
\caption{High-Feedback Target: Preliminary Holdout Results}
\label{tab:highfeedback}
\centering
\setlength{\tabcolsep}{3.3pt}
\begin{tabular}{llcccc}
\toprule
Input & Stage & AUROC & AUPRC & Brier & MLPD \\
\midrule
Measured & Prior & 0.8429 & 0.8794 & 0.2109 & $-0.6104$ \\
Measured & Final & 0.8460 & 0.8864 & 0.1901 & $-0.5606$ \\
Reconstructed & Prior & 0.8409 & 0.8807 & 0.2101 & $-0.6091$ \\
Reconstructed & Final & 0.8560 & 0.8936 & 0.1834 & $-0.5461$ \\
\bottomrule
\end{tabular}
\end{table}

The recurrent virtual sensor attains median held-out waveform correlations of 0.992, 0.995, and
0.989 for bounce, roll, and pitch rate, respectively. Its bounce-feature correlations range from
0.859 for derivative absolute peak to 0.992 for absolute impulse. The reconstructed preference
variant is not worse in this single split and slightly improves all four final metrics, but the
difference is too small and too dependent on one target to support a superiority claim.

\begin{figure*}[t]
\centering
\resultgraphic{true/plots/population/spike_slab_inclusion.png}{0.74\textwidth}
\caption{Posterior inclusion probabilities from the high-feedback executable snapshot. The plot
summarizes population feature support and is a diagnostic output rather than confirmatory evidence.}
\label{fig:diagnostics}
\end{figure*}

\draftnote{Regenerate the target-adaptation uncertainty plot with anonymous evaluator identifiers,
publication-size typography, and fixed feedback checkpoints before including it in the final paper.}

\subsection{Sparse-Feedback Targets}

Table~\ref{tab:sparse} shows the five sparse targets using measured channels. The final update uses
only 5--17 context labels, and holdouts contain as few as five episodes. Macro-average AUROC rises
from 0.6863 to 0.7008 and Brier score improves from 0.2441 to 0.2045. However, target S3 becomes
worse in AUROC after adaptation, and the near-perfect estimates for S4 or S5 are based on very few
positive--negative pairs. The result therefore supports neither universal improvement nor a precise
estimate of the mean effect. It instead demonstrates the regime in which posterior uncertainty and
proper probability scores are essential.

\begin{table}[t]
\caption{Sparse-Feedback Measured-Channel Results}
\label{tab:sparse}
\centering
\setlength{\tabcolsep}{4pt}
\begin{tabular}{lrrcccc}
\toprule
Target & $n_{\rm ctx}$ & $n_{\rm hold}$ & AUROC$_0$ & AUROC$_T$ & Brier$_0$ & Brier$_T$ \\
\midrule
S1 & 17 & 17 & 0.7083 & 0.7361 & 0.2115 & 0.2032 \\
S2 & 15 & 15 & 0.6481 & 0.7222 & 0.2304 & 0.2474 \\
S3 & 17 & 18 & 0.7000 & 0.4625 & 0.2404 & 0.2520 \\
S4 & 5  & 5  & 1.0000 & 0.8333 & 0.1880 & 0.1842 \\
S5 & 9  & 9  & 0.3750 & 0.7500 & 0.3500 & 0.1356 \\
\midrule
Macro & -- & -- & 0.6863 & 0.7008 & 0.2441 & 0.2045 \\
\bottomrule
\end{tabular}
\end{table}

For the reconstructed variant, macro Brier score similarly improves from 0.2438 to 0.2069 and mean
log predictive density from $-0.6823$ to $-0.5988$, while macro AUROC falls from 0.6986 to 0.6368.
This mixed result is consistent with the intended interpretation of reconstruction as a robustness
stress test rather than a performance contribution.

\subsection{Population Feature Structure}

Table~\ref{tab:feature-results} presents representative measured-channel posterior summaries from
the high-feedback split. Vertical acceleration peak-to-peak and longitudinal weighted RMS are the
first two features selected by the two predictive diagnostics, although their coefficient signs
are conditional on the correlated standardized design. The projection search requires seven
features plus the intercept to recover 95.4\% of the full-model KL reduction. The forward
conditional-information analysis retains six features before the next gain becomes nonpositive.

\begin{table*}[t]
\caption{Representative Population Feature Results in the High-Feedback Split}
\label{tab:feature-results}
\centering
\setlength{\tabcolsep}{5pt}
\begin{tabular}{lrrrrl}
\toprule
Feature & PIP & Mean coefficient & 95\% interval & Forward CMI (nat) & Diagnostic role \\
\midrule
Vertical acceleration peak-to-peak & 1.000 & $-0.694$ & $[-1.317,-0.104]$ & 0.07148 & First in projection and CMI \\
Longitudinal acceleration $W_d$ RMS & 1.000 & $-0.647$ & $[-1.407,0.094]$ & 0.01449 & Second in CMI; fifth in projection \\
Longitudinal acceleration standard deviation & 1.000 & $0.564$ & $[-0.094,1.189]$ & 0.00487 & Third in CMI; sixth in projection \\
Bounce-rate absolute impulse & 1.000 & $0.600$ & $[0.105,1.094]$ & 0.00318 & Early projection representative \\
Bounce-rate derivative $W_k$ RMS & 0.758 & $-0.155$ & $[-0.831,0.372]$ & 0.00525 & Later nonredundant CMI gain \\
Vertical acceleration vibration dose & 1.000 & $-0.070$ & $[-0.514,0.387]$ & -- & Completes 95\% projected subset \\
\bottomrule
\end{tabular}
\end{table*}

The posterior mean signs must not be read causally. Road severity, speed, vehicle setting, bump
identity, and evaluator history may be confounded, and correlated transient statistics can reverse
conditional signs. The defensible result is that several features jointly preserve predictive
structure while evaluator-specific coefficients vary; causal suspension-design statements require
controlled interventions.

\section{Discussion}
\label{sec:discussion}

\subsection{What Personalization Changes}

The high-feedback split suggests that the population prior already ranks episodes reasonably well,
whereas target feedback sharpens probability calibration and predictive density. This is a useful
distinction for intelligent-vehicle applications. A controller or recommendation layer often needs
a probability that reflects confidence, not only a ranking. The hierarchical prior avoids fitting a
new model from scratch, and the posterior trajectory makes it possible to define a future stopping
rule based on uncertainty rather than selecting the context with best holdout AUROC.

The sparse-target results show the other side of partial pooling. A few observations can move the
posterior in a helpful direction, but they can also be unrepresentative. Strong shrinkage toward the
population is desirable until the target supplies both good and bad examples. Future deployment
should therefore expose an abstention or ``request more feedback'' state and should not treat the
oracle peak as an online early-stopping mechanism.

\subsection{Interpretability and Sensor Robustness}

The three feature analyses are complementary. Spike-and-slab inclusion asks whether a feature is
active across posterior draws. Projection asks how compactly the full fitted predictions can be
reproduced. Cross-fitted conditional information asks whether a candidate improves held-out log
score after conditioning on the already selected features and evaluator identity. Agreement on
vertical peak-to-peak and longitudinal weighted RMS makes them stable descriptive candidates, but
not necessarily physical causes of good or bad judgments.

The reconstruction study indicates that a high-fidelity virtual bounce-rate signal can preserve
downstream probability quality in one held-out target. A rigorous missing-sensor treatment would
propagate reconstruction uncertainty into the preference likelihood rather than substituting a
point waveform. That extension is intentionally separated from the paper's present contribution.

\subsection{Limitations}

The number of evaluators is small relative to the number of episodes, and episode-level sample size
does not replace evaluator-level replication. Several sparse holdouts contain too few examples of
one class for stable AUROC inference. The current split does not explicitly group repeated passes by
bump, location, session, vehicle, or suspension setting, so within-group dependence and distribution
shift may be understated. Good/bad feedback is intentionally simple but cannot be equated with a
validated multi-attribute ride-comfort scale. The feature weights are vehicle-response proxies rather
than calibrated whole-body exposure measurements. Finally, the current artifact uses one population
chain and one principal split, and the available comparison set is not yet sufficient for a journal
claim of superiority.

\section{Confirmatory Study Required for Submission}
\label{sec:confirmatory}

The final T-IV evaluation should freeze preprocessing, priors, inclusion level, and stopping rules
before inspecting test results. At minimum, it should include the following.

\begin{enumerate}
  \item Leave one sufficiently sampled evaluator out in turn, fitting every normalization,
  feature-selection step, reconstruction model, and population posterior without that evaluator.

  \item For sparse feedback, subsample context budgets such as 0, 2, 5, 10, and 20 with repeated
  chronological starting points or session-aware blocks; keep a fixed group-held-out test set.

  \item Compare against a population-only logistic model, regularized independent models,
  mixed-effects or empirical-Bayes logistic regression, and at least one nonlinear sequence or
  tree-based baseline under the identical outer split.

  \item Report evaluator-macro proper scores with evaluator-level confidence intervals, calibration
  curves, sensitivity to class imbalance, prior-predictive checks, and multiple-chain convergence
  diagnostics. Treat episode-level intervals as secondary.

  \item Repeat the measured-versus-reconstructed comparison across outer folds and propagate or
  bound reconstruction error; do not tune the reconstruction network on final target evaluators.
\end{enumerate}

\draftnote{Replace all preliminary tables and the abstract's numerical claims with results from the
frozen confirmatory pipeline. Record software versions, hardware, run identifiers, wall-clock
costs, and inference seeds in a reproducibility appendix.}

\section{Conclusion}
\label{sec:conclusion}

This paper develops a hierarchical Bayesian framework for predicting evaluator-specific good/bad
ride preference after road bumps. The population posterior supplies a principled cold start, and
P\'{o}lya--Gamma personalization updates both predictions and uncertainty as sparse target feedback
arrives. Exact-zero inclusion, predictive projection, and conditional-information analysis retain
the interpretability of motion-derived features. Current executable results show improved
probability scores for one high-feedback target and on average across five sparse targets, while
also revealing substantial evaluator-level uncertainty and occasional negative adaptation. A
driver-held-out recurrent reconstruction preserves performance in the high-feedback snapshot but
remains a secondary robustness result. The next step is not a broader claim from the present split;
it is the confirmatory evaluator- and session-grouped evaluation specified above.

\section*{Acknowledgment}
\draftnote{Add funding, data-collection acknowledgments, conflicts of interest, and any required
ethics/consent statement after institutional confirmation.}

\bibliographystyle{IEEEtran}
\bibliography{refs}

% T-IV regular papers require short biographies in the final manuscript.
\ifdraftnotes
\vspace{1ex}
\noindent\draftnote{Add IEEE-format author biographies and photographs for the final files.}
\fi

\end{document}
